\section*{Bab \ref{ch:b2:integration} --- Pengintegralan}

\begin{solution}{exo:b2:integration:1}
$\int_0^1 \frac{\dd t}{\sqrt{t(1-t)}}$: di dekat $0$ ia $\sim t^{-1/2}$
(dengan $\alpha = \frac12 < 1$, jadi \omterm{def:b2:integration:improper}{konvergen}); di dekat $1$ ia $\sim
(1-t)^{-1/2}$, jadi \omterm{def:b2:integration:improper}{konvergen}. Karenanya \omterm{def:b2:integration:improper}{konvergen} (nilainya $\pi$, lewat
penyulihan $t = \sin^2\theta$).

$\int_1^\infty \frac{\ln t}{t^2}$: di sini $\frac{\ln t}{t^2} =
o(t^{-3/2})$, jadi \omterm{def:b2:integration:improper}{konvergen} (nilainya $1$ lewat pengintegralan parsial).

$\int_0^\infty \frac{\dd t}{1 + t^2\sin^2 t}$: divergen. Di dekat $t =
n\pi$, tulis $t = n\pi + u$: maka $\sin^2 t = \sin^2 u \leq u^2$, jadi
pada $\abs u \leq \frac{1}{n}$ berlaku $1 + t^2\sin^2 t \leq 1 + (n\pi +
1)^2u^2 \leq C n^2 u^2 + 1$; karenanya
\[
\int_{n\pi - 1/n}^{n\pi + 1/n} \frac{\dd t}{1 + t^2\sin^2 t}
\geq \int_{-1/n}^{1/n} \frac{\dd u}{1 + Cn^2u^2}
= \frac{2\arctan\sqrt C}{\sqrt C}\cdot\frac{1}{n} ,
\]
yakni suku sebuah deret bertipe harmonik yang divergen: jadi setelah
dijumlahkan atas $n$, antiturunannya tak terbatas.
\end{solution}

\begin{solution}{exo:b2:integration:2}
Sulihkan $u = \lambda t$:
\[
\int_0^\infty t^n \eu^{-\lambda t}\dd t
= \frac{1}{\lambda^{n+1}}\int_0^\infty u^n\eu^{-u}\dd u
= \frac{\Gamma(n+1)}{\lambda^{n+1}} = \frac{n!}{\lambda^{n+1}} .
\]
Dengan $t = \eu^{-u}$ (sehingga $\dd t = -\eu^{-u}\dd u$):
\[
\int_0^1 (\ln t)^n \dd t = \int_0^{\infty} (-u)^n \eu^{-u}\,\dd u
= (-1)^n\, n! .
\]
\end{solution}

\begin{solution}{exo:b2:integration:3}
Kekonvergenannya: integrannya $\leq \frac{1}{1+t^2}$ di dekat $\infty$
dan terbatas di dekat $0$ (sebab kedua faktornya terbatas di bawah,
menjauh dari $0$): jadi \omterm{def:b2:integration:improper}{konvergen} \omterm{def:b2:series:def}{mutlak} untuk setiap $x$. Dengan
menyulihkan $t = \frac1u$ (sehingga $\dd t = -\frac{\dd u}{u^2}$):
\[
I(x) = \int_0^\infty \frac{1}{\bigl(1 +
\frac1{u^2}\bigr)\bigl(1 + u^{-x}\bigr)}\cdot\frac{\dd u}{u^2}
= \int_0^\infty \frac{u^x}{(1 + u^2)(1 + u^x)}\,\dd u .
\]
Setelah kedua ungkapan $I(x)$ dijumlahkan:
\[
2I(x) = \int_0^\infty \frac{1 + t^x}{(1+t^2)(1+t^x)}\dd t
= \int_0^\infty \frac{\dd t}{1 + t^2} = \frac{\pi}{2} :
\]
jadi $I(x) = \frac\pi4$, yang tak bergantung pada $x$.
\end{solution}

\begin{solution}{exo:b2:integration:4}
Di dekat $0^+$, dengan $u = \abs{\ln t} \to \infty$. Jika $\alpha < 1$:
\omterm{def:b2:integration:improper}{konvergen} berapa pun $\beta$ (bandingkan dengan $t^{-\alpha'}$ untuk
$\alpha < \alpha' < 1$: sebab faktor logaritmanya terkalahkan). Jika
$\alpha > 1$: divergen berapa pun $\beta$ (bandingkan dengan
$t^{-\alpha''}$ dengan $1 < \alpha'' < \alpha$). Jika $\alpha = 1$:
sulihkan $t = \eu^{-u}$:
\[
\int_0^{1/2} \frac{\dd t}{t\,\abs{\ln t}^\beta}
= \int_{\ln 2}^{\infty} \frac{\dd u}{u^\beta},
\]
yang \omterm{def:b2:integration:improper}{konvergen} bila dan hanya bila $\beta > 1$. Ringkasnya: \omterm{def:b2:integration:improper}{konvergen}
bila dan hanya bila $\alpha < 1$, atau ($\alpha = 1$ dan $\beta > 1$)
--- yakni cerminan deret Bertrand.
\end{solution}

\begin{solution}{exo:b2:integration:5}
Kekontinuannya pada $\intco{0}{\infty}$: lewat dominasi
$\bigl|\frac{\eu^{-xt}}{1+t^2}\bigr| \leq \frac{1}{1+t^2}$, yang
terintegralkan dan seragam pada $x \geq 0$: jadi
\cref{thm:b2:integration:continuity}.

Kelas $C^2$ pada $\intoo{0}{\infty}$: pada $x \geq a > 0$, kedua turunan
pertamanya terhadap $x$, yakni $\frac{-t\,\eu^{-xt}}{1+t^2}$ dan
$\frac{t^2\eu^{-xt}}{1+t^2}$, terdominasi oleh $t\,\eu^{-at}$ dan
$\eu^{-at}$: jadi dua kali penerapan
\cref{thm:b2:integration:leibnizrule}. Maka
\[
F''(x) + F(x) = \int_0^\infty \frac{t^2 + 1}{1 + t^2}\,\eu^{-xt}\dd
t = \int_0^\infty \eu^{-xt}\dd t = \frac1x .
\]
Limitnya: $0 \leq F(x) \leq \int_0^\infty \eu^{-xt}\dd t = \frac1x \to
0$.
\end{solution}

\begin{solution}{exo:b2:integration:6}
Pada $\intcc{\varepsilon}{M}$, sulihkan $u = at$ dan $u = bt$ pada
kedua paruhnya:
\[
\int_\varepsilon^M \frac{f(at) - f(bt)}{t}\dd t
= \int_{a\varepsilon}^{aM}\frac{f(u)}{u}\dd u
- \int_{b\varepsilon}^{bM}\frac{f(u)}{u}\dd u
= \int_{a\varepsilon}^{b\varepsilon} \frac{f(u)}{u}\dd u
- \int_{aM}^{bM} \frac{f(u)}{u}\dd u .
\]
Potongan pertamanya: $f(u) = f(0) + o(1)$ di dekat $0$, dan
$\int_{a\varepsilon}^{b\varepsilon} \frac{\dd u}{u} = \ln\frac ba$: jadi
potongan itu menuju $f(0)\ln\frac ba$. Potongan keduanya: $f(u) \to
f(\infty)$, dengan perhitungan yang sama, jadi ia menuju
$f(\infty)\ln\frac ba$. Karenanya \omterm{def:b2:integration:improper}{integral tak wajarnya} \omterm{def:b2:integration:improper}{konvergen} ke
$\bigl(f(0) - f(\infty)\bigr)\ln\frac ba$.

Dengan $f(t) = \eu^{-t}$ (sehingga $f(0) = 1$, $f(\infty) = 0$), $a = 1$
dan $b = 2$:
\[
\int_0^\infty \frac{\eu^{-t} - \eu^{-2t}}{t}\dd t = \ln 2 .
\]
\end{solution}

\begin{solution}{exo:b2:integration:7}
Perluaslah integrannya oleh $0$ di luar $t = n$: $g_n(t) = (1 -
\frac tn)^n t^{x-1}\mathbf{1}_{t \leq n}$. Titik demi titik, $g_n(t) \to
\eu^{-t}t^{x-1}$ (yaitu limit bunga majemuk pada jilid Tahun ke-1).
Dominasinya: $\ln(1 - u) \leq -u$ memberi $(1 - \frac tn)^n \leq
\eu^{-t}$ pada $\intcc{0}{n}$, jadi $\abs{g_n(t)} \leq
\eu^{-t}t^{x-1} = \varphi(t)$, yang terintegralkan. Kekonvergenan
terdominasi memberi
\[
\int_0^n \Bigl(1 - \frac tn\Bigr)^n t^{x-1}\dd t
\xrightarrow[n\to\infty]{} \int_0^\infty \eu^{-t}t^{x-1}\dd t
= \Gamma(x) .
\]
(Menghitung ruas kirinya lewat pengintegralan parsial yang berulang memberi
bentuk hasil kali Euler $\Gamma(x) = \lim \frac{n!\,n^x}{x(x+1)\cdots(x+n)}$.)
\end{solution}

\begin{solution}{exo:b2:integration:8}
Fungsi $H$ dapat diturunkan terhadap $x$ (sebab integrannya $C^1$ pada
$x$, dengan turunan $-2x(1+t^2)\cdot\frac{\eu^{-x^2(1+t^2)}}{1+t^2} =
-2x\,\eu^{-x^2}\eu^{-x^2t^2}$, yang \omterm{def:b2:metric:continuity}{kontinu} dan terbatas pada himpunan
\omterm{def:b2:metric:compact}{kompak} dalam $x$, sedangkan dominasinya atas $t \in \intcc{0}{1}$ sepele):
\[
H'(x) = -2x\,\eu^{-x^2}\int_0^1 \eu^{-x^2t^2}\,\dd t
\overset{u = xt}{=} -2\,\eu^{-x^2}\int_0^x \eu^{-u^2}\,\dd u
= -G'(x),
\]
sebab $G'(x) = 2\eu^{-x^2}\int_0^x \eu^{-t^2}\dd t$ (lewat aturan rantai
pada kuadratnya dan teorema dasar kalkulus). Jadi $G + H$ bersifat
konstan, sama dengan $G(0) + H(0) = 0 + \int_0^1 \frac{\dd t}{1+t^2}
= \frac\pi4$.

Ketika $x \to \infty$: $0 \leq H(x) \leq \eu^{-x^2}\int_0^1 \dd t \to
0$, jadi $G(x) \to \frac\pi4$:
\[
\int_0^\infty \eu^{-t^2}\dd t = \sqrt{\frac\pi4} =
\frac{\sqrt\pi}{2} .
\]
(Akibatnya $\Gamma\bigl(\frac12\bigr) = 2\int_0^\infty
\eu^{-t^2}\dd t = \sqrt\pi$, lewat penyulihan $t = \sqrt u$.)
\end{solution}

\begin{solution}{exo:b2:integration:9}
Cauchy--Schwarz (jilid Tahun ke-1, yang sahih pada
$\intcc{\varepsilon}{M}$ lalu dilewatkan ke limitnya) yang diterapkan
pada pemfaktoran $t^{\frac{x+y}{2}-1}\eu^{-t} =
\bigl(t^{x-1}\eu^{-t}\bigr)^{1/2}
\bigl(t^{y-1}\eu^{-t}\bigr)^{1/2}$ memberi
\[
\Gamma\Bigl(\frac{x+y}{2}\Bigr) \leq
\Gamma(x)^{1/2}\,\Gamma(y)^{1/2}
\quad\Longrightarrow\quad
\ln\Gamma\Bigl(\frac{x+y}{2}\Bigr) \leq \frac{\ln\Gamma(x) +
\ln\Gamma(y)}{2} :
\]
jadi $\ln\Gamma$ cembung titik-tengah; dan karena ia \omterm{def:b2:metric:continuity}{kontinu}
(\cref{thm:b2:integration:gammaprops}), ia cembung
(\cref{exo:b2:realfun:8}). (Kelog-cembungan mengunci $\Gamma$
secara tunggal di antara semua interpolasi faktorialnya --- itulah
teorema Bohr--Mollerup, sebutir mutiara Tahun ke-3.)
\end{solution}

\begin{solution}{exo:b2:integration:10}
\begin{enumerate}
  \item Pada $x \geq a > 0$: turunan integrannya terhadap $x$ adalah
        $-\sin t\,\eu^{-xt}$, yang terdominasi oleh $\eu^{-at}$, jadi
        \cref{thm:b2:integration:leibnizrule} memberi $F'(x) =
        -\int_0^\infty \eu^{-xt}\sin t\,\dd t$. Dua kali pengintegralan
        parsial (atau lewat eksponensial kompleks) memberi
        \[
        \int_0^\infty \eu^{-xt}\sin t\,\dd t
        = \Im \int_0^\infty \eu^{(-x+\iu)t}\dd t
        = \Im\frac{1}{x - \iu} = \frac{1}{1 + x^2} .
        \]
  \item Berlaku $\abs{F(x)} \leq \int_0^\infty \eu^{-xt}\dd t = \frac1x
        \to 0$. Lalu mengintegralkan $F' = -\frac{1}{1+x^2}$ dari $x$
        sampai $\infty$ memberi $0 - F(x) = -\bigl(\frac\pi2 - \arctan
        x\bigr)$, jadi $F(x) = \frac\pi2 - \arctan x$.
  \item Dengan melewatkan $x \to 0^+$ beserta kekontinuan yang diterima
        tadi: $F(0^+) = \frac\pi2$, sedangkan $F(0) = \int_0^\infty
        \frac{\sin t}{t}\dd t$ (yakni integral Dirichlet yang
        \omterm{def:b2:integration:improper}{semi-konvergen}, \cref{ex:b2:integration:sint}): jadi nilainya
        $\frac\pi2$.
\end{enumerate}
\end{solution}

\begin{solution}{exo:b2:integration:11}
Kekonvergenannya: di dekat $0$ integrannya diperluas secara \omterm{def:b2:metric:continuity}{kontinu} oleh
nilai $1$ (sebab $\sin t \sim t$); di tak hingga ia $\leq t^{-2}$, jadi
\omterm{def:b2:integration:improper}{konvergen} \omterm{def:b2:series:def}{mutlak}. Pada $\intcc{\varepsilon}{M}$, integralkan secara
parsial dengan $u = \sin^2 t$ dan $v' = t^{-2}$:
\[
\int_\varepsilon^M \frac{\sin^2 t}{t^2}\dd t
= \Bigl[-\frac{\sin^2 t}{t}\Bigr]_\varepsilon^M
+ \int_\varepsilon^M \frac{2\sin t\cos t}{t}\dd t
= \Bigl[-\frac{\sin^2 t}{t}\Bigr]_\varepsilon^M
+ \int_{2\varepsilon}^{2M} \frac{\sin u}{u}\dd u
\]
(dengan $u = 2t$ pada integral terakhirnya). Kurungnya menuju $0$ di
kedua ujungnya (sebab $\sin^2\varepsilon/\varepsilon \leq \varepsilon$
dan $\sin^2 M/M \leq 1/M$), sedangkan integral terakhirnya menuju
$\int_0^\infty \frac{\sin u}{u}\dd u = \frac\pi2$
(\cref{exo:b2:integration:10}). Karenanya
\[
\int_0^\infty \Bigl(\frac{\sin t}{t}\Bigr)^{\!2}\dd t
= \frac\pi2 .
\]
\end{solution}

\begin{solution}{exo:b2:integration:12}
\begin{enumerate}
  \item Pengintegralan parsial dengan $u = \frac{-1}{2t}$ dan $v' =
        -2t\,\eu^{-t^2}$ (sehingga $v = \eu^{-t^2}$) memberi
        \[
        T(x) = \Bigl[\frac{-\eu^{-t^2}}{2t}\Bigr]_x^\infty
        - \int_x^\infty \frac{\eu^{-t^2}}{2t^2}\dd t
        = \frac{\eu^{-x^2}}{2x}
        - \int_x^\infty \frac{\eu^{-t^2}}{2t^2}\dd t .
        \]
        Alat yang sama pada integral barunya (dengan $u = \frac{-1}{4t^3}$
        dan $v' = -2t\,\eu^{-t^2}$):
        \[
        \int_x^\infty \frac{\eu^{-t^2}}{2t^2}\dd t
        = \frac{\eu^{-x^2}}{4x^3}
        - \frac34\int_x^\infty \frac{\eu^{-t^2}}{t^4}\dd t ,
        \]
        sehingga diperoleh kesamaan yang diumumkan tadi.
  \item Satu kali pengintegralan parsial lagi membatasi sisanya:
        \[
        \int_x^\infty \frac{\eu^{-t^2}}{t^4}\dd t
        = \frac{\eu^{-x^2}}{2x^5}
        - \frac52\int_x^\infty\frac{\eu^{-t^2}}{t^6}\dd t
        \leq \frac{\eu^{-x^2}}{2x^5},
        \]
        jadi $0 \leq \frac34\int_x^\infty t^{-4}\eu^{-t^2}\dd t
        \leq \frac{3}{8x^5}\eu^{-x^2}$. Membuang sisa yang positif
        pada kesamaan pertanyaan 1 memberi batas
        bawahnya; sedangkan membuang suku kedua yang negatif pada
        pengintegralan parsial pertamanya memberi $T(x) \leq
        \frac{\eu^{-x^2}}{2x}$. Setelah pengapitannya dibagi
        $\frac{\eu^{-x^2}}{2x}$: rasionya terapit
        di antara $1 - \frac{1}{2x^2}$ dan $1$, jadi $T(x)
        \sim \frac{\eu^{-x^2}}{2x}$.
  \item Mengiterasikan pengintegralan parsialnya menghasilkan deret formal
        \[
        T(x) \approx \frac{\eu^{-x^2}}{2x}\Bigl(1 -
        \frac{1}{2x^2} + \frac{1\cdot3}{(2x^2)^2} -
        \frac{1\cdot3\cdot5}{(2x^2)^3} + \cdots\Bigr),
        \]
        yang koefisien ke-$k$-nya $1\cdot3\cdots(2k-1) =
        \frac{(2k)!}{2^k k!}$ tumbuh lebih cepat daripada barisan
        geometri mana pun: sebab untuk $x$ yang tetap, suku
        $\frac{1\cdot3\cdots(2k-1)}{(2x^2)^k}$ menuju tak hingga
        (rasionya $\frac{2k+1}{2x^2} \to \infty$), sehingga
        deretnya divergen untuk setiap $x$. Ia sebuah uraian
        \emph{asimtotik}: dipotong pada orde tetap mana pun,
        galatnya seorde dengan suku pertama yang dibuang
        ketika $x \to \infty$ --- tetapi tak pernah menjadi deret
        yang \omterm{def:b2:integration:improper}{konvergen}. (Taksiran ekor ini adalah batas ekor Gauss
        yang baku pada bab peluang.)
\end{enumerate}
\end{solution}

\begin{solution}{pb:b2:integration:1}
\textbf{1.} Di $0^+$ integrannya $\sim t^{x-1}$: skala ujung
hingganya \omterm{def:b2:integration:improper}{konvergen} bila dan hanya bila $1 - x < 1$, yakni $x > 0$
(dan untuk $x \leq 0$, $t^{x-1} \geq t^{-1}$ divergen); sedangkan di
$+\infty$, $t^{x-1}\eu^{-t} = o(t^{-2})$ \omterm{def:b2:integration:improper}{konvergen} untuk setiap
$x$. Lalu $\Gamma(x) = \frac{\Gamma(x+1)}{x}$ dan $\Gamma(x+1)
\to \Gamma(1) = 1$ ketika $x \to 0^+$ (lewat kekontinuan,
\cref{thm:b2:integration:gammaprops}): jadi $\Gamma(x) \sim
\frac1x$.

\textbf{2.} Dengan $t = u^2$ dan $\dd t = 2u\,\dd u$:
\[
\Gamma\Bigl(\frac12\Bigr) = \int_0^\infty
t^{-1/2}\eu^{-t}\dd t
= \int_0^\infty \frac{\eu^{-u^2}}{u}\,2u\,\dd u
= 2\int_0^\infty \eu^{-u^2}\dd u = \sqrt\pi
\]
menurut \cref{exo:b2:integration:8}. Dengan $u = v/\sqrt2$:
\[
\int_\R \eu^{-v^2/2}\dd v
= 2\sqrt2\int_0^\infty \eu^{-u^2}\dd u
= \sqrt2\,\sqrt\pi = \sqrt{2\pi} .
\]

\textbf{3.} Benar untuk $n = 0$ (sebab kedua ruasnya $\sqrt\pi$). Jika
$\Gamma(n + \frac12) = \frac{(2n)!}{4^n n!}\sqrt\pi$, maka
persamaan fungsionalnya memberi
\[
\Gamma\Bigl(n + 1 + \frac12\Bigr)
= \Bigl(n + \frac12\Bigr)\Gamma\Bigl(n + \frac12\Bigr)
= \frac{2n+1}{2}\cdot\frac{(2n)!}{4^n n!}\sqrt\pi
= \frac{(2n+2)!}{4^{n+1}(n+1)!}\sqrt\pi ,
\]
dengan langkah terakhirnya berlaku karena $\frac{(2n+2)!}{(2n)!} =
(2n+2)(2n+1)$ dan $\frac{2n+1}{2} = \frac{(2n+2)(2n+1)}{4(n+1)}$.

\textbf{4.} \cref{thm:b2:integration:gammaprops} memberi
$\Gamma''(x) = \int_0^\infty t^{x-1}\eu^{-t}(\ln t)^2\dd t$
(lewat dua kali penerapan aturan Leibniz, dengan dominasi seperti pada
bukti teoremanya); integrannya $\geq 0$ dan tidak nol secara
identik, jadi $\Gamma'' > 0$: sehingga $\Gamma$ cembung tegas dan
$\Gamma'$ naik tegas. Karena $\Gamma(1) = \Gamma(2) = 1$, Rolle
menyediakan $x_0 \in \intoo12$ dengan $\Gamma'(x_0) = 0$; lalu
kemonotonan tegas $\Gamma'$ menjadikan $x_0$ satu-satunya nolnya,
dengan $\Gamma' < 0$ sebelumnya dan $\Gamma' > 0$ sesudahnya: jadi
$\Gamma$ turun pada $\intoo0{x_0}$, naik pada
$\intoo{x_0}\infty$, dan $x_0$ menjadi minimum tunggalnya.

\textbf{5.} Misalkan $k \in \N$ dan $x \geq 3$; pilihlah bilangan bulat
$n$ dengan $n + 1 \leq x < n + 2$ (sehingga $n \geq 1$). Menurut
kemonotonan pertanyaan 4 (yang sahih mulai $x_0 < 2$): $\Gamma(x)
\geq \Gamma(n + 1) = n!$, sedangkan $x^k \leq (n+2)^k$. Karena
$\frac{n!}{(n+2)^k} \to \infty$ (sebab faktorial mengalahkan pangkat,
jilid Tahun ke-1), berlaku $\frac{\Gamma(x)}{x^k} \geq
\frac{n!}{(n+2)^k} \to \infty$ ketika $x \to \infty$: jadi $x^k =
o(\Gamma(x))$.

\textbf{6.} Di dekat $0$ integrannya $\sim t^{x-1}$
(yang \omterm{def:b2:integration:improper}{konvergen} bila dan hanya bila $x > 0$), di dekat $1$ ia $\sim
(1-t)^{y-1}$ (bila dan hanya bila $y > 0$); kedua perbandingannya antara
fungsi positif, jadi $B(x,y)$ \omterm{def:b2:integration:improper}{konvergen} persis untuk $x, y > 0$.
Penyulihan $t \mapsto 1 - t$ menukar kedua faktornya: jadi $B(x,y) = B(y,x)$.

\textbf{7.} Di sini $B(x,1) = \int_0^1 t^{x-1}\dd t = \frac1x$.
Pengintegralan parsial pada $\intcc\varepsilon{1-\varepsilon}$ dengan $u
= (1-t)^y$ dan $v = \frac{t^x}{x}$ memberi
\[
\int t^{x-1}(1-t)^{y}\dd t
= \Bigl[\frac{t^x(1-t)^y}{x}\Bigr]
+ \frac{y}{x}\int t^{x}(1-t)^{y-1}\dd t ;
\]
dengan kurungnya nol di kedua ujungnya ketika $\varepsilon \to 0$ (sebab
$x > 0$ di $0$ dan $y > 0$ di $1$), sehingga tersisa $B(x, y+1) = \frac
yx\,B(x+1, y)$.

\textbf{8.} Karena $t + (1-t) = 1$:
\[
t^{x-1}(1-t)^{y-1} = t^{x}(1-t)^{y-1} + t^{x-1}(1-t)^{y},
\]
jadi $B(x,y) = B(x+1,y) + B(x,y+1)$. Pertanyaan 7 berbunyi $B(x+1,y) =
\frac xy B(x,y+1)$; setelah disulihkan,
\[
B(x,y) = \Bigl(\frac xy + 1\Bigr)B(x,y+1)
= \frac{x+y}{y}\,B(x,y+1),
\]
yakni $B(x,y+1) = \frac{y}{x+y}B(x,y)$; sedangkan relasi kembarannya
menyusul dari kesetangkupan pertanyaan 6.

\textbf{9.} Induksi pada $n$ dengan $m$ tetap: $B(m,1) = \frac1m =
\frac{(m-1)!\,0!}{m!}$, dan jika rumusnya berlaku di $n$, maka
\[
B(m, n+1) = \frac{n}{m+n}\,B(m,n)
= \frac{n}{m+n}\cdot\frac{(m-1)!(n-1)!}{(m+n-1)!}
= \frac{(m-1)!\,n!}{(m+n)!} .
\]
Setelah ditulis ulang: $B(m,n) = \frac{(m-1)!(n-1)!}{(m+n-1)!} =
\bigl[(m+n-1)\binom{m+n-2}{m-1}\bigr]^{-1}$.

\textbf{10.} Kedua ruas $B(x,n) =
\frac{\Gamma(x)\Gamma(n)}{\Gamma(x+n)}$ sama dengan $\frac1x$ di $n =
1$ (sebab $\Gamma(1) = 1$ dan $\Gamma(x+1) = x\Gamma(x)$). Jika keduanya
sepakat di $n$, maka lewat relasi penurunannya dan persamaan
fungsionalnya:
\[
B(x, n+1) = \frac{n}{x+n}\,B(x,n), \qquad
\frac{\Gamma(x)\Gamma(n+1)}{\Gamma(x+n+1)}
= \frac{n}{x+n}\cdot
\frac{\Gamma(x)\Gamma(n)}{\Gamma(x+n)} :
\]
kedua barisannya menuruti rekursi yang sama dari benih yang sama,
sehingga keduanya sepakat untuk setiap $n \in \N^*$ dan setiap $x > 0$.

\textbf{11.} Dengan $t = \sin^2\theta$ ($\theta \in
\intoo0{\pi/2}$, $\dd t = 2\sin\theta\cos\theta\,\dd\theta$), berlaku
$t^{x-1} = \sin^{2x-2}\theta$ dan $(1-t)^{y-1} =
\cos^{2y-2}\theta$, sehingga
\[
B(x,y) = \int_0^{\pi/2}\sin^{2x-2}\theta\,\cos^{2y-2}\theta
\cdot 2\sin\theta\cos\theta\,\dd\theta
= 2\int_0^{\pi/2}\sin^{2x-1}\theta\,\cos^{2y-1}\theta\,
\dd\theta .
\]

\textbf{12.} Ambil $y = \frac12$ (yang mematikan faktor kosinusnya) dan
$2x - 1 = n$: maka $B\bigl(\frac{n+1}2, \frac12\bigr) = 2W_n$, yakni
$W_n = \frac12 B\bigl(\frac{n+1}2,\frac12\bigr)$. Relasi penurunan pada
peubah pertamanya memberi
\[
\frac{W_n}{W_{n-2}}
= \frac{B\bigl(\frac{n-1}2 + 1, \frac12\bigr)}
{B\bigl(\frac{n-1}2, \frac12\bigr)}
= \frac{\frac{n-1}2}{\frac{n-1}2 + \frac12}
= \frac{n-1}{n} :
\]
yakni rekurensi Wallis, kali ini tanpa pengintegralan parsial atas
sinusnya --- sebab Bagian II sudah mengerjakannya sekali untuk selamanya.

\textbf{13.} Di sini $B\bigl(\frac12,\frac12\bigr) = 2W_0 =
2\cdot\frac\pi2 = \pi$, sedangkan
$\Gamma\bigl(\frac12\bigr)^2/\Gamma(1) = (\sqrt\pi)^2 = \pi$: jadi
rumus Euler berlaku di $\bigl(\frac12,\frac12\bigr)$.

\textbf{14.} Mengiterasikan $W_{2n} = \frac{2n-1}{2n}W_{2n-2}$ dari
$W_0 = \frac\pi2$ memberi
\[
W_{2n} = \frac\pi2\prod_{k=1}^{n}\frac{2k-1}{2k}
= \frac\pi2\cdot\frac{(2n)!}{4^n(n!)^2},
\]
sebab $\prod(2k-1) = \frac{(2n)!}{2^n n!}$ dan $\prod 2k = 2^n
n!$. Karenanya, dengan memakai pertanyaan 3:
\[
B\Bigl(n+\frac12, \frac12\Bigr) = 2W_{2n}
= \pi\,\frac{(2n)!}{4^n(n!)^2}
= \frac{(2n)!\sqrt\pi}{4^n n!}\cdot\frac{\sqrt\pi}{n!}
= \frac{\Gamma\bigl(n+\frac12\bigr)\Gamma\bigl(\frac12\bigr)}
{\Gamma(n+1)} .
\]
Kini tetapkan $x \in \frac12\N^*$. Rumus Euler berlaku di $(x,
\frac12)$: untuk $x$ bulat inilah pertanyaan 10 (dengan kesetangkupannya),
sedangkan untuk $x = n + \frac12$ inilah ungkapan di atas. Kedua ruas
rumus Euler menuruti rekursi penurunan $y \mapsto y + 1$
(pertanyaan 8 di kiri, persamaan fungsionalnya di kanan,
seperti pada pertanyaan 10): jadi induksi merambatkan rumusnya dari $y =
\frac12$ dan $y = 1$ ke setiap $y \in \frac12\N^*$. Karenanya rumus
Euler berlaku setiap kali $2x, 2y \in \N^*$.

\textbf{15.} Dengan $u = \frac{t}{1-t}$, yakni $t =
\frac{u}{1+u}$, $1 - t = \frac{1}{1+u}$, dan $\dd t =
\frac{\dd u}{(1+u)^2}$:
\[
B(x,y) = \int_0^\infty
\Bigl(\frac{u}{1+u}\Bigr)^{x-1}
\Bigl(\frac{1}{1+u}\Bigr)^{y-1}
\frac{\dd u}{(1+u)^2}
= \int_0^\infty \frac{u^{x-1}}{(1+u)^{x+y}}\,\dd u .
\]
Di $x = y = \frac12$, dengan $u = v^2$:
\[
\int_0^\infty \frac{u^{-1/2}}{1+u}\dd u
= \int_0^\infty \frac{2\,\dd v}{1+v^2} = \pi
= B\Bigl(\frac12,\frac12\Bigr) . \checkmark
\]

\textbf{16.} Satu kali pengintegralan parsial, untuk $1 \leq k \leq n$
dan $s > 0$ (dengan $u = (1 - t/n)^k$, $v = t^s/s$; dan suku batasnya
lenyap), memberi
\[
\int_0^n \Bigl(1-\frac tn\Bigr)^{\!k} t^{s-1}\dd t
= \frac{k}{ns}\int_0^n \Bigl(1-\frac tn\Bigr)^{\!k-1}
t^{s}\dd t .
\]
Berawal dari $k = n$, $s = x$ lalu mengiterasikannya $n$ kali:
\[
\int_0^n \Bigl(1-\frac tn\Bigr)^{\!n} t^{x-1}\dd t
= \frac{n(n-1)\cdots1}{n^n\,x(x+1)\cdots(x+n-1)}
\int_0^n t^{x+n-1}\dd t
= \frac{n!}{n^n}\cdot
\frac{n^{x+n}}{x(x+1)\cdots(x+n)} ,
\]
yakni $\dfrac{n!\,n^x}{x(x+1)\cdots(x+n)}$.

\textbf{17.} Menurut \cref{exo:b2:integration:7}, ruas kirinya menuju
$\Gamma(x)$ (lewat kekonvergenan terdominasi dengan pendominasi
$t^{x-1}\eu^{-t}$); sedangkan ruas kanannya adalah kuosien Gauss:
\[
\Gamma(x) = \lim_{n\to\infty}
\frac{n!\,n^x}{x(x+1)\cdots(x+n)} .
\]

\textbf{18.} Setelah logaritma kuosien $G_n(x)$ pada pertanyaan 16
diambil lalu $\ln(x+k) = \ln k + \ln(1 + x/k)$ dipecah untuk $k
\geq 1$:
\[
\ln G_n(x) = x\ln n - \ln x - \sum_{k=1}^n
\ln\Bigl(1+\frac xk\Bigr)
= -\ln x + x(\ln n - H_n)
+ \sum_{k=1}^n\Bigl(\frac xk -
\ln\Bigl(1+\frac xk\Bigr)\Bigr).
\]
Untuk $u \geq 0$ berlaku $u - \frac{u^2}2 \leq \ln(1+u) \leq u$, jadi
suku umumnya berada di $\intcc{0}{x^2/(2k^2)}$: sehingga deretnya
\omterm{def:b2:integration:improper}{konvergen} (lewat pembandingan dengan $\sum k^{-2}$). Karena $\ln n - H_n
\to -\gamma$ (\cref{ex:b2:comparison:harmonic}) dan $\ln G_n(x)
\to \ln\Gamma(x)$ (pertanyaan 17 dan kekontinuan $\ln$):
\[
\ln\Gamma(x) = -\ln x - \gamma x
+ \sum_{k=1}^\infty\Bigl(\frac xk -
\ln\Bigl(1+\frac xk\Bigr)\Bigr) .
\]

\textbf{19.} Di $x = \frac12$, penyebutnya adalah
$\prod_{k=0}^n\bigl(k+\frac12\bigr) =
\frac{(2n+1)!}{2^{2n+1}n!}$ (kalikan saja setengahnya), sehingga
\[
G_n\Bigl(\frac12\Bigr)
= \frac{n!\,\sqrt n\;2^{2n+1}n!}{(2n+1)!}
= \frac{2\sqrt n\;4^n}{(2n+1)\binom{2n}{n}} .
\]
Dengan $\binom{2n}{n} \sim \frac{4^n}{\sqrt{\pi n}}$
(\cref{ex:b2:comparison:centralbinomial}):
\[
G_n\Bigl(\frac12\Bigr)
\sim \frac{2\sqrt n\,\sqrt{\pi n}}{2n+1}
\longrightarrow \sqrt\pi = \Gamma\Bigl(\frac12\Bigr) .
\]
Jadi $\sqrt\pi$ milik koefisien binomial pusat (yang berasal dari
Wallis, sehingga dari konstanta Stirling) dan $\sqrt\pi$ milik integral
Gauss adalah bilangan yang sama.

\textbf{20.} Kesamaannya murni aljabar: kalikan
$\frac{n!\,n^x}{x(x+1)\cdots(x+n)}$ dengan $\frac{nx}{x+n+1}$ lalu
serap $x$ ke dalam hasil kalinya dan $n$ ke dalam $n^x$. Dengan
melewatkan $n \to \infty$: ruas kirinya menuju $\Gamma(x+1)$ (yakni
Gauss di $x+1$), sedangkan ruas kanannya menuju $\Gamma(x)\cdot x\cdot
1$ karena $\frac{n}{x+n+1} \to 1$: jadi $\Gamma(x+1) = x\Gamma(x)$ ---
dipulihkan tanpa satu pun pengintegralan parsial.

\textbf{21.} Dengan $u = t^a$, $t = u^{1/a}$, dan $\dd t = \frac1a
u^{1/a - 1}\dd u$:
\[
\int_0^\infty \eu^{-t^a}\dd t
= \frac1a\int_0^\infty u^{\frac1a - 1}\eu^{-u}\dd u
= \frac1a\,\Gamma\Bigl(\frac1a\Bigr)
= \Gamma\Bigl(1 + \frac1a\Bigr) .
\]
Ketika $a \to +\infty$ (sepanjang barisan mana pun): $\eu^{-t^a} \to 1$
untuk $0 < t < 1$, $\to \eu^{-1}$ di $t = 1$, dan $\to 0$ untuk $t > 1$;
sedangkan untuk $a \geq 2$ dominasilah oleh $\mathbf 1_{t \leq 1} +
\eu^{-t^2}\mathbf 1_{t > 1}$ (sebab $t^a \geq t^2$ untuk $t \geq 1$),
yang terintegralkan. Kekonvergenan terdominasi: jadi integralnya menuju
$\int_0^1 1\,\dd t = 1$ --- sebagaimana mestinya, sebab $\Gamma(1 +
\frac1a) \to \Gamma(1) = 1$ berkat kekontinuannya.

\textbf{22.} Dengan $u = t^n$ dan $\dd t = \frac1n u^{1/n - 1}\dd
u$:
\[
\int_0^1 \frac{\dd t}{\sqrt{1-t^n}}
= \frac1n\int_0^1 u^{\frac1n-1}(1-u)^{-1/2}\dd u
= \frac1n\,B\Bigl(\frac1n, \frac12\Bigr) .
\]
Untuk $n = 1$: $B\bigl(1,\frac12\bigr) = B\bigl(\frac12,1\bigr) = 2$,
yang cocok dengan $\int_0^1\frac{\dd t}{\sqrt{1-t}} = 2$. Untuk $n = 2$:
$\frac12 B\bigl(\frac12,\frac12\bigr) = \frac\pi2 = \arcsin 1$.
Sedangkan untuk $n = 4$, nilai $\frac14 B\bigl(\frac14,\frac12\bigr)$
adalah konstanta lemniskat: tanpa bentuk tertutup yang elementer.

\textbf{23.} Dengan mengiterasikan persamaan fungsionalnya:
\[
\frac{1}{\Gamma(x)}\int_0^\infty t^{x+k-1}\eu^{-t}\dd t
= \frac{\Gamma(x+k)}{\Gamma(x)}
= (x+k-1)(x+k-2)\cdots x ,
\]
yakni faktorial naik dengan $k$ faktor. Di $x = 1$:
$\Gamma(1+k)/\Gamma(1) = k!$, yakni momen $\eu^{-t}$ pada
\cref{exo:b2:integration:2}.

\textbf{24.} Sulihkan $t = \frac{1+s}2$ (sehingga $s \in
\intoo{-1}1$, $\dd t = \frac{\dd s}2$, dan $t(1-t) =
\frac{1-s^2}4$):
\[
B(x,x) = \int_{-1}^{1}\Bigl(\frac{1-s^2}{4}\Bigr)^{x-1}
\frac{\dd s}{2}
= 4^{1-x}\int_0^1 (1-s^2)^{x-1}\dd s
\]
(sebab integrannya genap). Lalu $s = \sqrt v$ (sehingga $\dd s =
\frac{\dd v}{2\sqrt v}$):
\[
B(x,x) = \frac{4^{1-x}}{2}\int_0^1
v^{-1/2}(1-v)^{x-1}\dd v
= 2^{1-2x}\,B\Bigl(\frac12, x\Bigr) .
\]
Untuk $2x \in \N^*$ setiap argumen yang tampak berada di
$\frac12\N^*$, jadi rumus Euler (pertanyaan 14) berlaku bagi kedua
ruasnya:
\[
\frac{\Gamma(x)^2}{\Gamma(2x)}
= 2^{1-2x}\,
\frac{\Gamma\bigl(\frac12\bigr)\Gamma(x)}
{\Gamma\bigl(x+\frac12\bigr)}
\quad\Longleftrightarrow\quad
\Gamma(x)\,\Gamma\Bigl(x+\frac12\Bigr)
= 2^{1-2x}\sqrt\pi\;\Gamma(2x) .
\]
Pemeriksaan langsung di $x = n$: ruas kirinya $(n-1)!\cdot
\frac{(2n)!\sqrt\pi}{4^n n!} = \frac{(2n)!\sqrt\pi}{4^n n}$, sedangkan
ruas kanannya $2\cdot4^{-n}\sqrt\pi\,(2n-1)! =
\frac{(2n)!\sqrt\pi}{4^n n}$: jadi sama.

\textbf{25.} (i) Pengintegralan parsial menghasilkan $B(x,y+1) =
\frac yx B(x+1,y)$, yakni satu-satunya kesamaan yang darinya mengalir
setiap relasi penurunan, nilai bulat dan setengah bulatnya, serta
rekurensi Wallis. (ii) Kekonvergenan terdominasi mengubah
integral elementer $\int_0^n(1-t/n)^n t^{x-1}$ menjadi
$\Gamma(x)$ (yakni rumus Gauss, pertanyaan 17) lalu menghitung
limit $a \to \infty$ pada pertanyaan 21. (iii) Dari bab
perbandingan kita mengimpor konstanta Euler (yakni $\ln n - H_n \to
-\gamma$, pertanyaan 18) dan asimtotik koefisien binomial pusat
(pertanyaan 19) --- yaitu rumus Stirling yang menyamar. (iv)
Rumus Euler $B(x,y) = \Gamma(x)\Gamma(y)/\Gamma(x+y)$ kini
terbukti untuk $y \in \N^*$ dengan $x > 0$ sembarang (pertanyaan
10) dan untuk semua pasangan setengah bulat (pertanyaan 14);
sedangkan kasus umum $x, y > 0$ menunggu perhitungan integral ganda
pada bab integral lipat, yang memfaktorkan
$\Gamma(x)\Gamma(y)$ atas seperempat bidang.

\end{solution}
